\section{Discussion}

In this section, we reflect on our research, discuss lessons learned and the system's generalizability, followed by limitations and future work. 

\subsection{Lessons Learned}
\textbf{Benefits of structured strategy representation for human-LLMs co-design process.} Our evaluation results indicate that having a structured representation for coordination strategy is essential for the experience of the design process. The structure helps align human and LLMs' intentions, enforcing both sides to think and communicate based on the same set of concepts and abstraction levels, which effectively reduces mutual misunderstandings caused by the ambiguity of natural language. Moreover, a fixed structure enables the specialized design of visual organization and enhancement, facilitating strategy comprehension for humans, allowing them to efficiently explore more strategy possibilities generated by LLMs. Additionally, structured representation fosters a structured exploration process, making users inclined towards a more systematic design process for coordination strategies.

\textbf{Visualizing LLMs' prior knowledge for agent coordination.} While many works directly leverage LLMs' prior knowledge to generate coordination strategy, how to effectively extract and visualize this prior knowledge to aid strategy design has not been explored. We find that compared to just getting a single answer from LLMs, users prefer to have a more panoramic understanding of LLM's prior knowledge when they want to optimize a certain part of the strategy. For example, instead of just letting LLM select some agents that might contribute the a certain task, visualizing how each agent scores for the capabilities LLMs deem important using an interactive heat map would be more insightful and systematic for strategy design. 

\subsection{System Generalizability}
Besides coordinating agents to collaboratively solve goal-oriented tasks, our system has the potential to be generalized to various applications. For example, users can use our system to coordinate agents to simulate and analyze human activities (e.g. playing Werewolf games, competing in debate tournaments, conducting multi-party negotiations) that involve multiple people. Those simulations are not only interesting for entertainment but also valuable resources for studying human/AI society and evaluating specific capabilities of LLMs. In the future, we plan to extend the structure representation of \textit{AgentCoord} to support more social interaction types for more flexible simulation purposes. We also plan to allow users to upload their customized interaction types.



\subsection{Limitations and Future Work}
\textit{AgentCoord} currently only supports coordinating agents to collaborate in a plain text environment, which contains text-form key objects. An interesting future work is to support multi-modal key objects in the environment and allow agents with multi-model capabilities to collaborate. For example, a writer agent generates a story, and an illustrator agent creates corresponding visuals.

\textit{AgentCoord} currently only supports static coordination strategy design. In some scenarios, users may want to dynamically adjust the strategy during collaboration execution to adapt to circumstances. For example, the user may coordinate multiple agents to help conduct literature research: when an agent finds a paper that interests the user, the user may want to adjust the plan to allocate a group of agents to read and discuss it, and allocate another group of agents to investigate the background of its authors. Future research can be conducted to help users conduct dynamic coordination strategy design.






